% Part: incompleteness
% Chapter: theories-computability
% Section: computably-axiomatizable

\documentclass[../../../include/open-logic-section]{subfiles}

\begin{document}

\olfileid{inc}{tcp}{cax}

\olsection{\printtoken{S}{axiomatizable} تيورۍ}

يوه تيوري~$\Th{T}$ هله~\emph{!!{axiomatizable}} بلل کېږي چې د بديهي
اصولو يو محاسبه کېدونکے سټ~$A$ ولري۔ (دا ويل چې~$A$ د~$\Th{T}$
د بديهي اصولو سټ دے، معنا ئې دا ده چې
$T = \Setabs{!B}{A \Proves !B}$۔)
\textbf{د ژباړې د نسخې سمون (OLCMP-093):} سرچينې په دې سټ-جوړوونکې
بڼه کښې يوه نښه هم د بديهي اصولو د سټ او هم د اړوندې جملې دپاره
کارولې وه؛ دلته جمله په بله نښه بېله شوې ده، څو تعريف هماغه تيوري
وټاکي چې له ټاکلي بديهي سټ څخه د اشتقاق وړ ده۔
د طبيعي عددونو هره «معقوله» بديهي کول به دا خاصيت لري۔ په ځانګړې
توګه، هره هغه تيوري چې د بديهي اصولو متناهي سټ ولري،
!!{axiomatizable} وي۔

\begin{lem}
فرض کړئ~$\Th{T}$، !!{axiomatizable} ده۔ بيا~$\Th{T}$،
!!{computably
enumerable} ده۔
\end{lem}

\begin{proof}
فرض کړئ~$A$ د~$\Th{T}$ د بديهي اصولو يو محاسبه کېدونکے سټ دے۔ د
دې د ټاکلو دپاره چې~$!A \in T$ ده که نۀ، يوازې د بديهي اصولو څخه
د~$!A$ د !!a{derivation} پلټنه وکړئ۔

لږ په بله وينا، $!A$ هله او يوازې هله په~$\Th{T}$ کښې ده چې
په~$A$ کښې د بديهي اصولو يو متناهي لست~$!B_1$، \dots، $!B_k$ او
په لومړۍ درجې منطق کښې د~$(!B_1 \land \dots \land !B_k) \lif !A$
يو !!a{derivation} وي۔ خو موږ لا پوهېږو چې هر هغه سټ چې تعريف ئې د
«داسې \dots شته چې \dots» په بڼه وي، !!{c.e.} وي، په دې شرط چې
دويم «\dots» محاسبه کېدونکے وي۔
\end{proof}

\end{document}
